\section*{Overview}

A 2D Fenwick tree is the direct generalization of the ordinary Fenwick tree from arrays to grids. It is the cleanest
tool for:

\begin{itemize}
  \item point updates on a matrix,
  \item prefix-sum queries on rectangles,
  \item rectangle sums via inclusion-exclusion.
\end{itemize}

\section*{When to Use It}

Use it when:

\begin{itemize}
  \item the data lives on a 2D grid,
  \item updates touch one cell at a time,
  \item queries ask for sums over axis-aligned rectangles.
\end{itemize}

If the grid is huge but only a few coordinates ever appear, combine it with coordinate compression first.

\section*{Core Idea}

The one-dimensional Fenwick tree decomposes a prefix into power-of-two blocks. The 2D version does the same in both
coordinates. Updating one cell climbs through all ancestors in the \(x\) dimension and, inside each of those, through
all ancestors in the \(y\) dimension.

\section*{Key Insight}

The logarithms multiply, not the dimensions of the whole grid. That is why the cost is \(O(\log n \log m)\) rather
than linear in one direction.

\section*{Worked Problem}

\paragraph{Problem.}
Maintain a board where stars can be added or removed from cells, and answer the number of stars inside any rectangle
\([x_1, x_2] \times [y_1, y_2]\).

\paragraph{Why 2D Fenwick fits.}
Rectangle sums reduce to four prefix sums:
\[
\texttt{sum}(x_2,y_2)-\texttt{sum}(x_1-1,y_2)-\texttt{sum}(x_2,y_1-1)+\texttt{sum}(x_1-1,y_1-1).
\]

\section*{Correctness Intuition}

Each cell contributes to exactly the 2D Fenwick blocks that cover it, just as in one dimension. Prefix queries collect
exactly the blocks needed to tile the requested prefix rectangle without overlap.

\section*{Complexity Analysis}

\begin{itemize}
  \item point update: \(O(\log n \log m)\),
  \item prefix sum: \(O(\log n \log m)\),
  \item rectangle sum: \(O(\log n \log m)\).
\end{itemize}

\section*{Implementation}

The sample code supports point add and rectangle sum queries on a 1-indexed grid.

\section*{Common Pitfalls}

\begin{itemize}
  \item Forgetting 1-indexing.
  \item Accidentally allocating a full \(n \times m\) structure when compression was needed.
  \item Using it when the problem really needs rectangle updates or range updates, which require extra machinery.
\end{itemize}

\section*{Variants / Extensions}

\begin{itemize}
  \item Coordinate-compressed sparse 2D Fenwick tree.
  \item Higher-dimensional Fenwick trees.
  \item 2D segment tree when the update/query model is more complicated.
\end{itemize}

\section*{Practice Problems}

\begin{itemize}
  \item Point update and rectangle sum on a grid.
  \item Dynamic counting of marked cells inside rectangles.
  \item Compressed-coordinate rectangle query problems.
\end{itemize}

\section*{References}

\begin{itemize}
  \item \href{https://cp-algorithms.com/data_structures/fenwick.html}{cp-algorithms: Fenwick Tree}.
  \item \href{https://usaco.guide/plat/2DRQ}{USACO Guide: 2D Range Queries}.
  \item \href{https://codeforces.com/catalog}{Codeforces Catalog}.
\end{itemize}
